\section{Preliminaries}\label{sec:preliminaries}

This section reviews the graphon framework and the background results used
later in the paper.  We recall the relevant entropy-maximization theory of
Kenyon, Radin, Ren, and Sadun, state the two-sided parametrization derived
from their results in Appendix~\ref{app:krrs-analytic-extension}.

\subsection{Graphons}\label{sec:graphons}

A graphon is a symmetric measurable function
$W:[0,1]^2\to[0,1]$.  We think of $W(x,y)$ as the probability of an edge
between the continuum vertices $x$ and $y$, so that $W$ specifies a
random graph model on the vertex space $[0,1]$.  We write $\W_0$ for the
graphon space.

The \emph{cut norm} of a symmetric measurable
$F:[0,1]^2\to\RR$ is
\[
  \|F\|_\cut
  =\sup_{S,T}\left|\int_{S\times T}F(x,y)\dd x\dd y\right|,
\]
the supremum being over measurable $S,T\subseteq[0,1]$, and the \emph{cut
distance} between two graphons is
\[
  \delta_\cut(U,W)=\inf_\sigma\|U-W^\sigma\|_\cut,
  \qquad
  W^\sigma(x,y):=W(\sigma x,\sigma y),
\]
where $\sigma$ ranges over the measure-preserving a.e. bijections of $[0,1]$.
We say that two graphons $U$ and $W$ are \emph{weakly isomorphic} if
$\delta_\cut(U,W)=0$, and write $\widetilde{\W}_0$ for the quotient of
$\W_0$ by this equivalence relation.  The cut distance $\delta_\cut$ can be lifted to a metric
on $\widetilde{\W}_0$, and the resulting metric space $(\widetilde{\W}_0,\delta_\cut)$ is compact.  When
discussing uniqueness or compactness, we always work in
$\widetilde{\W}_0$; ``unique up to relabelling'' means unique as an element
of $\widetilde{\W}_0$, and in particular two graphons differing by a
measure-preserving relabelling of $[0,1]$ are identified.  Every homomorphism density
$t(H,\cdot)$ is $\delta_\cut$-continuous, and $I_p$ is $\delta_\cut$-lower
semicontinuous.  See~\cite{LovaszGraphLimits} for all of this.  For a graph $G$,
we write $\delta_\cut(G,W)=\delta_\cut(W_G,W)$, where $W_G$ is the step graphon
associated with $G$.\footnote{Namely, partition $[0,1]$ into $|V(G)|$
intervals of equal length, indexed by the vertices of $G$, and set $W_G=1$
on $I_i\times I_j$ when $ij\in E(G)$, and $W_G=0$ otherwise.  A different
ordering of the vertices changes $W_G$ only by a measure-preserving
relabelling, so $\delta_\cut(G,W)$ is well defined.}

We say that a graphon $W$ is \emph{bipodal} if there is a measurable set
$A\subset[0,1]$ such that, writing $A^c=[0,1]\setminus A$, the graphon $W$
is almost everywhere equal to a constant on each of
\[
  A\times A,
  \qquad
  (A\times A^c)\cup(A^c\times A),
  \qquad
  A^c\times A^c.
\]
Equivalently, after relabelling, a bipodal graphon is described by parameters
$(q_{11},q_{12},q_{22},c)$, where $c=|A|$ denotes the Lebesgue measure, or block
size, of $A$.  The parameters $q_{11}$ and $q_{22}$ are the edge densities
within $A$ and $A^c$, respectively, and $q_{12}$ is the edge density between
the two blocks.  Following the terminology in the entropy-maximization
literature, we also call the two blocks the two \emph{podes} of the graphon.

For a finite graph $H$, its homomorphism density in $W$ is
\[
  t(H,W)
  =
  \int_{[0,1]^{|V(H)|}}\prod_{ij\in E(H)} W(x_i,x_j)
  \prod_{i\in V(H)} dx_i.
\]
In particular, the edge density is
\[
  e(W)=t(K_2,W)=\int_{[0,1]^2} W(x,y)\,dx\,dy.
\]
The constant graphon with value $u$ is denoted by $W\equiv u$.

When $H$ is $d$-regular, we use the generalized H\"older inequality in the
following form, valid for every graphon $W$:
\begin{equation}\label{eq:generalized-holder}
  t(H,W)\le
  \left(\int_{[0,1]^2}W(x,y)^d\,dx\,dy\right)^{|E(H)|/d}.
\end{equation}



\subsection{Entropy maximization problem}
The entropy-maximization problem reviewed here is closely related to the
Chatterjee--Varadhan variational problem~\eqref{eq:graphon-variational-problem}.
The connection between such constrained entropy problems and large
deviations for the uniform random graph $G(n,M)$ was established by Dembo
and Lubetzky~\cite{DemboLubetzky2018}.  For related work on constrained
graphon entropy maximization, see
\cite{AlvaradoRamonWang2026,KenyonRadinRenSadun2017,
KenyonRadinRenSadun2017b,NeemanRadinSadun2023}.
Define the graphon entropy functional by
\[
  s(W)=-\frac12\int_{[0,1]^2}\left[W(x,y)\log W(x,y)+(1-W(x,y))\log(1-W(x,y))\right]dx\,dy.
\]
For fixed $p$ and fixed edge density, minimizing $I_p$ is equivalent to
maximizing $s$, since
\begin{equation}\label{eq:relative-entropy-identity}
  I_p(W)=-2s(W)-\log(1-p)+e(W)\log\frac{1-p}{p}.
\end{equation}

We say that a graph $H$ is \emph{$d$-starlike} if it has a vertex of degree
$d$ and all other vertices have degree either $1$ or $d$.  In particular,
every $d$-regular graph is $d$-starlike.  Fix such a graph $H$ and write
$m:=|E(H)|$.  For $0<\eps<1$ and $\eps^m<\tau<1$, consider the
entropy-maximization problem with edge density $\eps$ and $H$-density $\tau$:
\[
  \sup\{s(W):e(W)=\eps,\ t(H,W)=\tau\}.
\]
The constant graphon $W\equiv\eps$ has $H$-density $\eps^m$.
Accordingly, define
\[
  \vartheta:=\tau-\eps^m
\]
which is the positive $H$-density surplus above the constant-graphon value.

For the small-positive-surplus regime considered here, Kenyon, Radin, Ren,
and Sadun~\cite[Theorem~1.1]{kenyon2014entropy} proved
that for every $\eps\ne(d-1)/d$ there is a threshold
$\tau_0(\eps)>\eps^m$ such that, whenever
$\eps^m<\tau<\tau_0(\eps)$, this problem has a unique maximizer up to
relabelling, and the maximizer is bipodal.  Its four bipodal parameters are
real-analytic in $(\eps,\tau)$ throughout this positive-surplus region, which
is therefore open.  Consequently, for every fixed
$\eps_0\ne(d-1)/d$, one can choose a neighbourhood $U$ of $\eps_0$ and a
single $\Delta>0$ such that the KRR--S conclusions hold for every $\eps\in U$
whenever $0<\vartheta<\Delta$.

In addition to establishing this analytic parametrization for
$\vartheta>0$, KRR--S determined its one-sided behaviour as
$\vartheta\downarrow0$.  To state their result for the limiting cross density
$q_{12}$, write
\[
  S_0(u):=-\frac12\bigl[u\log u+(1-u)\log(1-u)\bigr]
\]
for the one-edge entropy density, so that
$s(W)=\int_{[0,1]^2}S_0(W(x,y))\,dx\,dy$, and define
\[
  \psi_d(\eps,z):=
  \frac{2\bigl[S_0(z)-S_0(\eps)-S_0'(\eps)(z-\eps)\bigr]}
       {z^d-\eps^d-d\eps^{d-1}(z-\eps)},
\]
where the removable singularity at $z=\eps$ is filled in continuously.  For
each $\eps\in(0,1)$, the function $z\mapsto\psi_d(\eps,z)$ has a unique
maximizer on $(0,1)$, denoted by $\zeta_d(\eps)$; see
\cite[Theorem~3.3 and Section~4]{kenyon2014entropy}.  With this notation, for
each fixed $\eps\ne(d-1)/d$, their result gives
\[
  q_{22}\longrightarrow\eps,
  \qquad
  q_{12}\longrightarrow\zeta_d(\eps),
  \qquad
  c=O(\vartheta)
  \qquad(\vartheta\downarrow0).
\]
For later reference, when $d=2$, a direct computation using
$S_0(1-\eps)=S_0(\eps)$ and
$S_0'(1-\eps)=-S_0'(\eps)$ gives $\zeta_2(\eps)=1-\eps$.
\Cref{thm:krrs-analytic-extension} below also implies that $\zeta_d$ is real-analytic on
$(0,1)\setminus\{(d-1)/d\}$.

Because the four parameters depend on $(\eps,\tau)$, we call the function
\[
  (\eps,\tau)\longmapsto(q_{11},q_{12},q_{22},c)
\]
the KRR--S parameter map.  Equivalently, using
$\tau=\eps^m+\vartheta$, we view it as a function of
$(\eps,\vartheta)$.  Our later analysis differentiates its components and
uses their Taylor expansions at the boundary $\vartheta=0$.  For this
purpose, the parameter map must be defined on an open neighbourhood of the
boundary, whereas KRR--S define it only for $\vartheta>0$.

The following theorem serves two purposes.  It records the KRR--S
positive-surplus conclusions in the form used throughout the rest of the
paper, so that later arguments can refer directly to this theorem.  It also
adds the two-sided real-analytic continuation of the parameter map through
$\vartheta=0$ that those arguments require.
\begin{theorem}[Two-sided extension of the KRR--S parametrization]\label{thm:krrs-analytic-extension}
Let $H$ be a $d$-starlike graph with $m=|E(H)|\ge2$.
Fix an edge density
\[
  \eps_0\in(0,1)\setminus\left\{\frac{d-1}{d}\right\}.
\]
Then, there are an open neighbourhood
$U\subseteq(0,1)\setminus\{(d-1)/d\}$ of $\eps_0$ and a constant $\Delta>0$
with the following property.  For every $\eps\in U$ and every $\tau$ such that
$0<\vartheta:=\tau-\eps^m<\Delta$, the entropy maximizer among graphons satisfying
\[
  e(W)=\eps,
  \qquad
  t(H,W)=\tau
\]
is unique up to relabelling and is bipodal.  After relabelling so that the
smaller block is the first block, the corresponding bipodal parameters
\[
  (q_{11},q_{12},q_{22},c)
\]
define a parameter map on $U\times(0,\Delta)$ that extends real-analytically
to $U\times(-\Delta,\Delta)$.  On $U\times\{0\}$, the extension satisfies
\[
  c(\eps,0)=0,
  \qquad
  q_{22}(\eps,0)=\eps,
  \qquad
  q_{12}(\eps,0)=\zeta_d(\eps)\ne\eps.
\]
Moreover,
\[
  c(\eps,\vartheta)=O(\vartheta)
  \qquad(\vartheta\downarrow0),
\]
uniformly for $\eps\in U$.
\end{theorem}
At $\vartheta=0$, the first pode has size zero, and hence the bipodal graphon
reduces to the constant graphon $W\equiv\eps$.  Consequently, the values of
$q_{11}$ and $q_{12}$ on the null first pode are not determined by the graphon
itself; their values at $\vartheta=0$ are the distinguished ones selected by
analytic continuation from $\vartheta>0$.  Likewise, for $\vartheta<0$, the
extended parameter functions are only analytic continuations and do not
describe entropy maximizers.

The analytic continuation in \Cref{thm:krrs-analytic-extension} is constructed in
Appendix~\ref{app:krrs-analytic-extension} by two applications of the analytic implicit
function theorem, using Theorems~1.1 and~3.3 of KRR--S.